\subsection{结合律}

以下引理补充 Ch. V, §8, No.~5 的命题 11：\(^{2}\)

引理 1. --- 对 \(1 \leqslant i \leqslant n\)，令 \(\mathrm{X}_i, \mathrm{Y}_i\) 为两个局部紧空间，\(\mu_i\) 为 \(\mathrm{X}_i\) 上的测度，\(\varphi_i\) 为从 \(\mathrm{X}_i\) 到 \(\mathrm{Y}_i\) 的连续映射。令 \(\mathrm{X} = \prod_i \mathrm{X}_i\)、\(\mathrm{Y} = \prod_i \mathrm{Y}_i\)、\(\mu = \bigotimes_i \mu_i\)，并令 \(\varphi\) 为从 \(\mathrm{X}\) 到 \(\mathrm{Y}\) 的映射，即各 \(\varphi_i\) 的乘积。若 \(\varphi\) 是 \(\mu\)-适当的且 \(\mu_i \neq 0\) 对每个 \(i\) 都成立，则各 \(\varphi_i\) 是 \(\mu_i\)-适当的，并且 \(\varphi(\mu) = \bigotimes_i \varphi_i(\mu_i)\)。

不妨设 \(\mu_{i}\) 为正测度且 \(n = 2\)。取 \(f_{1} \in \mathcal{K}_{+}(Y_{1})\)。由于 \(\mu_{2} \neq 0\)，存在 \(f_{2} \in \mathcal{K}_{+}(Y_{2})\)，使得 \(f_{2} \circ \varphi_{2}\) 不是 \(\mu_{2}\)-可忽略的。函数 \((x_{1}, x_{2}) \mapsto f_{1}(\varphi_{1}(x_{1})) f_{2}(\varphi_{2}(x_{2}))\) 本质上 \(\mu\)-可积且连续，因而 \(\mu\)-可积。因此存在 \(x_{2} \in X_{2}\)，使得 \(f_{2}(\varphi_{2}(x_{2})) \neq 0\)，并且函数 \(x_{1} \mapsto f_{1}(\varphi_{1}(x_{1})) f_{2}(\varphi_{2}(x_{2}))\) 是 \(\mu_{1}\)-可积的。因此 \(f_{1} \circ \varphi_{1}\) 是 \(\mu_{1}\)-可积的，这就证明了 \(\varphi_{1}\) 是 \(\mu_{1}\)-适当的。对 \(\varphi_{2}\) 作同样论证。于是根据 Ch. V, §8, No.~5 的命题 11，\(\varphi(\mu) = \bigotimes_{i} \varphi_{i}(\mu_{i})\)。

以下引理补充 Ch. V, §6, No.~3 的命题 4。\(^{3}\)

引理 2. --- 设 \(\mathrm{T}, \mathrm{T}', \mathrm{T}''\) 为三个局部紧空间，\(\mu\) 为 \(\mathrm{T}\) 上的测度，\(\pi\) 为 \(\mu\)-可测的从 \(\mathrm{T}\) 到 \(\mathrm{T}'\) 的映射，\(\pi'\) 为从 \(\mathrm{T}'\) 到 \(\mathrm{T}''\) 的连续映射，且 \(\pi'' = \pi' \circ \pi\)。若 \(\pi''\) 是 \(\mu\)-适当的，则 \(\pi\) 是 \(\mu\)-适当的，\(\pi'\) 是 \(\pi(\mu)\)-适当的，并且 \(\pi''(\mu) = \pi' (\pi(\mu))\)。

令 \(K'\) 为 \(T'\) 的紧子集。则 \(K'' = \pi'(K')\) 是紧的，因而 \(\pi''(K'')\) 本质上 \(\mu\)-可积；于是 \(\overline{\pi}(K') \subset \overline{\pi''}(K'')\) 本质上 \(\mu\)-可积，从而 \(\pi\) 是 \(\mu\)-适当的。然后根据 Ch. V, §6, No.~3 的命题 4，\(\pi'\) 是 \(\pi(\mu)\)-适当的，且 \(\pi''(\mu) = \pi'(\pi(\mu))\)。

命题 1. --- 设 \(\mathrm{X}_{ij}\)（\(1 \leqslant i \leqslant m, 1 \leqslant j \leqslant n_i\)）、\(\mathrm{Y}_i\)（\(1 \leqslant i \leqslant m\)）和 \(\mathbf{Z}\) 为局部紧空间；对每个 \(i\)，令 \(\varphi_i\) 为从 \(\mathrm{X}_i = \prod_j \mathrm{X}_{ij}\) 到 \(\mathrm{Y}_i\) 的映射；令 \(\varphi\) 为各 \(\varphi_i\) 的乘积，即从 \(\mathrm{X} = \prod_i \mathrm{X}_i\) 到 \(\mathrm{Y} = \prod_j \mathrm{Y}_i\) 的映射；令 \(\psi\) 为从 \(\mathrm{Y}\) 到 \(\mathbf{Z}\) 的映射。

\begin{enumerate}
\def\labelenumi{(\roman{enumi})}
\tightlist
\item
  令 \(\mu_{ij}\) 分别为定义在 \(\mathrm{X}_{ij}\) 上的测度，满足：对每个 \(i\)，\(\mu_{ij}\)（\(1 \leqslant j \leqslant n_i\)）关于 \(\varphi_i\) 可卷积，并且测度 \(*|_{\mu_{ij}}|\) 关于 \(\psi\) 可卷积；则 \(\mu_{ij}\)，其中 \(1 \leqslant i \leqslant m\)、\(1 \leqslant j \leqslant n_i\)，关于 \((\psi \circ \varphi)\) 可卷积，并且
\end{enumerate}

\[
\underset {i, j} {\ast} \mu_ {i j} = \underset {i} {\ast} \left(\underset {j} {\ast} \mu_ {i j}\right).\tag{3}
\]

\begin{enumerate}
\def\labelenumi{(\roman{enumi})}
\setcounter{enumi}{1}
\tightlist
\item
  假设 \(\psi\) 和各 \(\varphi_{i}\) 连续，且令 \(\mu_{ij}\)（\(\neq 0\)）分别为定义在 \(X_{ij}\) 上的测度，并且关于 \((\psi \circ \varphi)\) 可卷积；则对每个 \(i\)，\(\mu_{ij}\)（\(1 \leqslant j \leqslant n_{i}\)）关于 \(\varphi_{i}\) 可卷积，测度 \(*_{j} |\mu_{ij}|\) 关于 \(\psi\) 可卷积，并且公式 (3) 成立。
\end{enumerate}

只需考虑所涉及的测度全都满足 \(\geqslant 0\) 的情形。

先置于 (i) 的假设下。映射 \(\varphi\) 对 \(\bigotimes_{i,j} \mu_{ij}\) 是适当的，并且

\[
\varphi \left(\bigotimes_ {i, j} \mu_ {i j}\right) = \bigotimes_ {i} \varphi_ {i} \left(\bigotimes_ {j} \mu_ {i j}\right) = \bigotimes_ {i} \left(* _ {j} \mu_ {i j}\right)
\]

（Ch. V, §8, No.~5, 命题 11）。映射 \(\psi \circ \varphi\) 对 \(\bigotimes_{i,j} \mu_{ij}\) 是适当的，并且

\[
(\psi \circ \varphi) \left(\bigotimes_ {i, j} \mu_ {i j}\right) = \psi \left(\bigotimes_ {i} \left(* _ {j} \mu_ {i j}\right)\right) = * _ {i} \left(* _ {j} \mu_ {i j}\right)
\]

（Ch. V, §6, 命题 4）。因此 \(\mu_{ij}\)（\(1 \leqslant i \leqslant m\)、\(1 \leqslant j \leqslant n_i\)）关于 \((\psi \circ \varphi)\) 可卷积，且公式 (3) 成立。

现在置于 (ii) 的假设下。首先，引理 2 证明 \(\varphi\) 对 \(\bigotimes_{i,j} \mu_{ij}\) 是适当的。于是引理 1 证明对每个 \(i\)，\(\varphi_i\) 对 \(\bigotimes_{i} \mu_{ij}\) 是适当的，并且

\[
\varphi \Big (\bigotimes_ {i, j} \mu_ {i j} \Big) = \bigotimes_ {i} \Big (* _ {j} \mu_ {i j} \Big).
\]

根据引理 2，\(\psi\) 对 \(\bigotimes_{i}\left(*_{j}\mu_{ij}\right)\) 是适当的。命题得证。

推论。--- 设 \(X_{i}, X_{i}^{\prime}\)（\(1 \leqslant i \leqslant n\)）、Y、\(Y^{\prime}\) 为局部紧空间；令 \(\varphi, \varphi^{\prime}\) 分别为从 \(X = \prod_{i} X_{i}\) 到 Y 以及从 \(X^{\prime} = \prod_{i} X_{i}^{\prime}\) 到 \(Y^{\prime}\) 的连续映射；令 \(f_{i}\) 为从 \(X_{i}\) 到 \(X_{i}^{\prime}\) 的连续映射（\(1 \leqslant i \leqslant n\)），g 为从 Y 到 \(Y^{\prime}\) 的连续映射，满足 \(\varphi^{\prime} \circ f = g \circ \varphi\)，其中 f 是从 X 到 \(X^{\prime}\) 的映射，即各 \(f_{i}\) 的乘积。令 \(\mu_{i}\) 分别为定义在 \(X_{i}\) 上的测度，且均满足 \(\neq 0\)。则以下两个断言等价：

\begin{enumerate}
\def\labelenumi{(\roman{enumi})}
\item
  \(f_{i}\) 是 \(\mu_{i}\)-适当的（对所有 \(i\)），并且测度 \(f_{i}(|\mu_{i}|)\) 关于 \(\varphi'\) 可卷积；
\item
  \(\mu_{i}\) 关于 \(\varphi\) 可卷积，并且 \(g\) 对 \(*_{\varphi}(|\mu_i|)\) 是适当的。
\end{enumerate}

此外，当这些断言成立时，

\[
\left. \ast_ {\varphi^ {\prime}} \left(f _ {i} \left(\mu_ {i}\right)\right) = g \left(\ast_ {\varphi} \left(\mu_ {i}\right)\right) = \ast_ {g \circ \varphi} \left(\mu_ {i}\right). \right.\tag{4}
\]

事实上，令 \(h = \varphi' \circ f = g \circ \varphi\)。根据命题 1，条件 (i) 和 (ii) 分别都等价于以下条件：

\begin{enumerate}
\def\labelenumi{(\roman{enumi})}
\setcounter{enumi}{2}
\tightlist
\item
  \(\mu_{i}\) 关于 \(h\) 可卷积。
\end{enumerate}

在此情形下，

\[
\left. \ast_ {\varphi^ {\prime}} \left(f _ {i} \left(\mu_ {i}\right)\right) = \ast_ {h} \left(\mu_ {i}\right) = g \left(\ast_ {\varphi} \left(\mu_ {i}\right)\right). \right.
\]

